\subsection{Metric variation and its normalization}
\label{sec:exact-arithmetic-metric-variation}

Write \(I(L_0,\ldots,L_n)\) for an arithmetic intersection satisfying
symmetry, tensor additivity and the divisor recursion.  Use metric
denominators \(h_L\), so that \(\|s\|=|s|/h_L\); the archimedean
term in the recursion is minus the integral of \(\log\|s\|\).
Denote its original mixed integral by
\(\langle u;L_1,\ldots,L_n\rangle_\infty\), retaining every complex
component and its weight.

For two smooth metrics \(A,B\) of the same multidegree, put
\(\delta=\log h_B-\log h_A\).  Then
\begin{equation}\label{eq:metric-variation-single-slot}
 I(L_1,\ldots,L_n,B)-I(L_1,\ldots,L_n,A)
       =\langle\delta;L_1,\ldots,L_n\rangle_\infty .
\end{equation}
Indeed, the trivial line with denominator \(e^\delta\) has section \(1\),
empty divisor and logarithmic section norm \(-\delta\).
The two minus signs in the recursion give the positive sign in
\eqref{eq:metric-variation-single-slot}.

There is a short way to justify locality even when homogeneous metric
data have arbitrary values at degenerate coordinate blocks.
Let \(\mathcal C\) be the set where every block is nonzero and use
\[
                    h_C(z)=\prod_k
                       \left(\sum_i|z_{ki}|^2\right)^{1/2}.
\]
This is a smooth positive metric on \(\mathcal C\), and vanishes
outside it.  If two same-degree metrics agree on \(\mathcal C\),
their tensor products with \(C\) agree everywhere; tensor additivity
and cancellation give equal intersections.  Extend \(\delta\) by zero
outside \(\mathcal C\) and apply this observation to
\(A\otimes(\mathcal O,e^\delta)\) and \(B\), proving
\eqref{eq:metric-variation-single-slot}.
If \(A,B\) agree only on the variety, \(\delta\) vanishes at every
point of its mixed integral, so \eqref{eq:metric-variation-single-slot}
also gives locality on the variety.
Thus changing a smooth representative does not change the intersection.

Replacing the factors one at a time now gives
\begin{equation}\label{eq:metric-variation-top-intersection}
 I(B^{n+1})-I(A^{n+1})
   =\sum_{j=0}^{n}
          \langle\delta;B^{[j]},A^{[n-j]}\rangle_\infty .
\end{equation}
Here brackets indicate repeated arguments.  If the original mixed
integrals are identified on a smooth compact complex model with the
corresponding Chern integrals, write
\(h_A=h_0e^\varphi\), \(h_B=h_0e^\psi\) and
\(\theta=dd^c\log h_0\), with \(dd^c=i\partial\bar\partial/\pi\).
Equation \eqref{eq:metric-variation-top-intersection} is precisely
\begin{equation}\label{eq:metric-variation-energy-factor}
 I(B^{n+1})-I(A^{n+1})
                =(n+1)\bigl(E_\theta(\psi)-E_\theta(\varphi)\bigr),
\end{equation}
with the original weighted sum over components understood.
The energy is unnormalized; the mixed difference formula is
\cite[Proposition~2.1]{boundary-bbgz}.  There is no further factor \(1/2\):
the half in the cusp comparison is the given difference
\(\psi-\varphi=U/2\).

For integers \(p\geq0\) and \(q\geq1\), multilinearity also gives
\(q^{-(n+1)}I((L^{\otimes q}\otimes A^{\otimes p})^{n+1})
 =I((L+(p/q)A)^{n+1})\), interpreted as the usual polynomial,
whose constant term is \(I(L^{n+1})\).
This establishes the algebraic normalization.  Identifying the original
mixed integrals and common vertical curvature with the specified compact
model remains a separate geometric requirement.
